\section{Local Limit and Modular Classification of the Infinite System}

In an infinite UHF algebra, there is in general no global trace-class
density matrix within the algebra. The correct object is the
net
\begin{equation}
\label{rm:eq:modular-local-net}
\HHM^{(m)}=-\log\rho_m
\end{equation}
with compatible-state restrictions, or the modular generator in the
GNS representation. The tracial expectation \(\mathrm{id}\otimes\tau_9\) does
not automatically preserve a nontracial Gibbs state and does not
replace state compatibility.

\begin{theorem}[Modular type under faithful persistence of both channels]
\label{rm:thm:modular-typeIII}
Suppose that:
\begin{enumerate}[label=(\roman*)]
\item each cut is a product of faithful qutrit states with weights
\(w_{90}\) and \(w_{120}\);
\item the inclusions preserve the state and the channel label;
\item infinitely many factors of each type occur, both with positive
asymptotic density;
\item no additional spectral quotients are introduced, and the
product representation is factorial.
\end{enumerate}
Then the additive group of logarithms of ratios of eigenvalues is
\begin{equation}
\label{rm:eq:modular-ratio-group}
90\Dmod\Z+120\Dmod\Z=30\Dmod\Z.
\end{equation}
The GNS closure ITPFI is the hyperinfinitesimal factor
\begin{equation}
\label{rm:eq:modular-typeIII-lambda}
\boxed{
III_{\lambda_\partial},
\qquad
\lambda_\partial=\ee^{-30\Dmod}
=0.9966696464689573786108299769\ldots .}
\end{equation}
\end{theorem}

\begin{proof}
At each link, the eigenvalue quotients are powers of \(\ee^{-w_m}\). Infinite persistence of both types generates the subgroup \(w_{90}\Z+w_{120}\Z\). Since \(\gcd(90,120)=30\), it follows that \cref{rm:eq:modular-ratio-group}. The Araki--Woods asymptotic ratio criterion for the faithful product yields the hyperfinite ITPFI factor with parameter \(\ee^{-30\Dmod}\) \cite{rm:ArakiWoods1968,rm:Takesaki}.
\end{proof}

The theorem states its actual hypotheses. If one of the channels disappears,
if the product is incompatible, or if the ratio group changes, the
type must be recalculated.
